% !PRINTED 20
\weberpage{20}
%<<CONT
Equality occurs here only when the quotient $z:c$ is real and
negative.

\begin{center}
\rule{0.18\linewidth}{0.4pt}
\end{center}

It remains to say a word about the most important aid of algebra,
literal calculation. This aid is so generally used that the term
literal calculation is sometimes employed outright as a synonym for
algebra. We presuppose the rules for calculating with such literal
expressions. Equations between literal expressions can be of two
kinds: either they are so-called identities, that is, the two
expressions set equal can be transformed by the rules of calculation
until they agree exactly. Such literal equations yield correct
numerical equations when the letters are replaced by any real or
complex numbers, provided that division by zero is not required. The
letters in such equations are often called variables, because one may
imagine, without ever reaching a contradiction, that different
numerical values are successively substituted for them.

Another kind of equation between literal expressions does not have the
character of an identity. It instead contains a requirement imposed on
the numbers that may be substituted for the letters without error.
Algebra has the task of determining numerical values that satisfy such
a requirement, that is, of solving the equation. In these equations
the letters are also called ``\gesperrt{unknowns}.'' Very often one
and the same equation contains letters of two kinds: some for which
arbitrary numerical values are to be substituted, and others whose
numerical value is to be determined.

\begin{center}
\rule{0.24\linewidth}{0.4pt}
\end{center}
